\documentclass[11pt,a4paper]{article}
\usepackage[utf8]{inputenc}
\usepackage[T1]{fontenc}
\usepackage[ngerman]{babel}
\usepackage{helvet}
\renewcommand{\familydefault}{\sfdefault}
\usepackage[left=2cm, right=2cm, top=2cm, bottom=2cm]{geometry}
\usepackage{amsmath,amssymb}
\usepackage{array}
\usepackage{booktabs}
\usepackage{tikz}
\usepackage{pgfplots}
\pgfplotsset{compat=1.18}
\usepackage{fancyhdr}
\usepackage{tcolorbox}
\usepackage{enumitem}

\pagestyle{fancy}
\fancyhf{}
\fancyhead[L]{\small Physik Klasse 10E}
\fancyhead[R]{\small Klausurlernheft Kinematik}
\fancyfoot[C]{\small Seite \thepage}
\renewcommand{\headrulewidth}{0.4pt}

\setlength{\parindent}{0pt}
\setlength{\fboxsep}{8pt}

% Abschnitt-Befehl im Layout B Stil
\newcounter{abschnitt}
\newcommand{\abschnitt}[1]{%
\stepcounter{abschnitt}%
\vspace{0.5cm}%
\noindent\fbox{\parbox{\dimexpr\textwidth-2\fboxsep-2\fboxrule}{%
\textbf{\theabschnitt. #1}%
}}%
\vspace{0.3cm}%
}

% Formelbox
\newcommand{\formelbox}[2]{%
\begin{center}
\fbox{\parbox{0.9\textwidth}{%
\textbf{#1}\\[0.2cm]
#2%
}}
\end{center}
}

% Beispiel-Umgebung
\newcounter{beispiel}
\newenvironment{beispiel}[1]{%
\stepcounter{beispiel}%
\vspace{0.3cm}%
\noindent\textbf{Beispiel \thebeispiel: #1}\\[0.2cm]
}{\vspace{0.3cm}}

% Aufgaben-Umgebung
\newcounter{uebung}
\newcommand{\uebung}[1]{%
\stepcounter{uebung}%
\vspace{0.2cm}%
\noindent\textbf{Übung \theuebung:} #1%
\vspace{0.2cm}%
}

\begin{document}

% === TITELSEITE ===
\begin{center}
\vspace*{2cm}

\noindent\fbox{\parbox{0.9\textwidth}{
\centering
\vspace{0.5cm}
{\Huge\bfseries Klausurlernheft}\\[0.5cm]
{\LARGE Kinematik}\\[0.3cm]
{\large Geradlinige Bewegungen}\\[0.5cm]
{\normalsize Klasse 10E Gymnasium}\\[0.3cm]
{\normalsize Schuljahr 2024/25}\\[0.5cm]
}}

\vspace{2cm}

\begin{tabular}{|l|p{8cm}|}
\hline
\textbf{Themen} & Gleichförmige Bewegung, beschleunigte Bewegung, freier Fall, senkrechter Wurf, Diagramme \\
\hline
\textbf{Hilfsmittel} & Tafelwerk (kein Taschenrechner!) \\
\hline
\textbf{Hinweis} & Alle Rechnungen sind kopfrechenbar mit $g = 10\,\mathrm{m/s^2}$ \\
\hline
\end{tabular}

\vspace{2cm}

\begin{tabular}{@{}p{4cm}p{8cm}@{}}
Name: & \hrulefill \\[0.5cm]
Klasse: & \hrulefill \\
\end{tabular}

\end{center}

\newpage

% === INHALTSVERZEICHNIS ===
\tableofcontents

\newpage

% === TEIL 1: FORMELSAMMLUNG ===
\section{Formelsammlung}

\abschnitt{Einheitenumrechnung}

\formelbox{Geschwindigkeit: km/h $\leftrightarrow$ m/s}{
\begin{center}
$v\,[\mathrm{m/s}] = \dfrac{v\,[\mathrm{km/h}]}{3{,}6}$ \qquad\qquad $v\,[\mathrm{km/h}] = v\,[\mathrm{m/s}] \cdot 3{,}6$
\end{center}
}

\textbf{Wichtige Umrechnungen (auswendig!):}
\begin{center}
\begin{tabular}{|c|c|c|c|c|c|}
\hline
\textbf{km/h} & 36 & 72 & 90 & 108 & 180 \\
\hline
\textbf{m/s} & 10 & 20 & 25 & 30 & 50 \\
\hline
\end{tabular}
\end{center}

\abschnitt{Gleichförmige Bewegung}

Bei gleichförmiger Bewegung ist die Geschwindigkeit $v$ konstant.

\formelbox{Weg-Zeit-Gesetz (gleichförmig)}{
\begin{center}
$s(t) = v \cdot t + s_0$
\end{center}
\begin{tabular}{@{}ll}
$s$ & zurückgelegter Weg in m \\
$v$ & Geschwindigkeit in m/s \\
$t$ & Zeit in s \\
$s_0$ & Startposition in m \\
\end{tabular}
}

\formelbox{Geschwindigkeit}{
\begin{center}
$v = \dfrac{s}{t}$ \qquad bzw. \qquad $v = \dfrac{\Delta s}{\Delta t}$
\end{center}
}

\abschnitt{Gleichmäßig beschleunigte Bewegung}

Bei beschleunigter Bewegung ändert sich die Geschwindigkeit gleichmäßig.

\formelbox{Geschwindigkeit-Zeit-Gesetz}{
\begin{center}
$v(t) = a \cdot t + v_0$
\end{center}
\begin{tabular}{@{}ll}
$v$ & Geschwindigkeit in m/s \\
$a$ & Beschleunigung in m/s² \\
$t$ & Zeit in s \\
$v_0$ & Anfangsgeschwindigkeit in m/s \\
\end{tabular}
}

\formelbox{Weg-Zeit-Gesetz (beschleunigt)}{
\begin{center}
$s(t) = \dfrac{1}{2} a \cdot t^2 + v_0 \cdot t + s_0$
\end{center}
}

\formelbox{Beschleunigung}{
\begin{center}
$a = \dfrac{\Delta v}{\Delta t} = \dfrac{v - v_0}{t}$
\end{center}
}

\formelbox{Zeitlose Formel (ohne $t$)}{
\begin{center}
$v^2 = v_0^2 + 2 \cdot a \cdot s$
\end{center}
}

\newpage

\abschnitt{Bremsvorgang}

Beim Bremsen ist $a < 0$ (Verzögerung).

\formelbox{Bremszeit}{
\begin{center}
$t = \dfrac{v_0}{|a|}$
\end{center}
}

\formelbox{Bremsweg}{
\begin{center}
$s = \dfrac{v_0^2}{2 \cdot |a|}$
\end{center}
Herleitung aus der zeitlosen Formel mit $v = 0$:
\begin{center}
$0 = v_0^2 - 2 \cdot |a| \cdot s \quad\Rightarrow\quad s = \dfrac{v_0^2}{2 \cdot |a|}$
\end{center}
}

\abschnitt{Freier Fall}

Spezialfall der beschleunigten Bewegung mit $a = g$ und $v_0 = 0$.

\formelbox{Freier Fall (Start aus der Ruhe)}{
\begin{center}
$s = \dfrac{1}{2} g \cdot t^2$ \qquad\qquad $v = g \cdot t$ \qquad\qquad $v = \sqrt{2 \cdot g \cdot s}$
\end{center}
\textbf{Wichtig:} $g \approx 10\,\mathrm{m/s^2}$ (kopfrechenbar!)
}

\formelbox{Fallzeit aus Höhe $h$}{
\begin{center}
$t = \sqrt{\dfrac{2h}{g}}$
\end{center}
}

\textbf{Kopfrechenbare Wurzeln:}
\begin{center}
\begin{tabular}{|c|c|c|c|c|c|c|c|c|c|}
\hline
$\sqrt{4}$ & $\sqrt{9}$ & $\sqrt{16}$ & $\sqrt{25}$ & $\sqrt{36}$ & $\sqrt{49}$ & $\sqrt{64}$ & $\sqrt{81}$ & $\sqrt{100}$ \\
\hline
2 & 3 & 4 & 5 & 6 & 7 & 8 & 9 & 10 \\
\hline
\end{tabular}
\end{center}

\abschnitt{Senkrechter Wurf nach oben}

Der Körper wird nach oben geworfen und von der Schwerkraft abgebremst.

\formelbox{Senkrechter Wurf nach oben}{
Geschwindigkeit: $v(t) = v_0 - g \cdot t$\\[0.2cm]
Höhe: $h(t) = v_0 \cdot t - \dfrac{1}{2} g \cdot t^2$\\[0.3cm]
Steigzeit (bis $v = 0$): $t_{\text{steig}} = \dfrac{v_0}{g}$\\[0.3cm]
Maximale Höhe: $h_{\text{max}} = \dfrac{v_0^2}{2g}$
}

\textbf{Wichtig:} Am höchsten Punkt ist $v = 0$. Die Rückfallgeschwindigkeit entspricht betragsmäßig der Abwurfgeschwindigkeit (Energieerhaltung).

\newpage

\abschnitt{Diagramme}

\formelbox{Bedeutung von Steigung und Fläche}{
\begin{center}
\begin{tabular}{|l|c|c|}
\hline
\textbf{Diagramm} & \textbf{Steigung bedeutet} & \textbf{Fläche bedeutet} \\
\hline
$s(t)$-Diagramm & Geschwindigkeit $v$ & --- \\
\hline
$v(t)$-Diagramm & Beschleunigung $a$ & Weg $s$ \\
\hline
\end{tabular}
\end{center}
}

\textbf{Typische $v(t)$-Diagramme:}

\begin{center}
\begin{tikzpicture}
\begin{axis}[
    width=5cm, height=3.5cm,
    xlabel={$t$},
    ylabel={$v$},
    xmin=0, xmax=5,
    ymin=0, ymax=4,
    axis lines=left,
    xtick=\empty,
    ytick=\empty,
    title={\small Gleichförmig}
]
\addplot[thick, black] coordinates {(0,2) (5,2)};
\end{axis}
\end{tikzpicture}
\hspace{0.5cm}
\begin{tikzpicture}
\begin{axis}[
    width=5cm, height=3.5cm,
    xlabel={$t$},
    ylabel={$v$},
    xmin=0, xmax=5,
    ymin=0, ymax=4,
    axis lines=left,
    xtick=\empty,
    ytick=\empty,
    title={\small Beschleunigt}
]
\addplot[thick, black] coordinates {(0,0) (5,3.5)};
\end{axis}
\end{tikzpicture}
\hspace{0.5cm}
\begin{tikzpicture}
\begin{axis}[
    width=5cm, height=3.5cm,
    xlabel={$t$},
    ylabel={$v$},
    xmin=0, xmax=5,
    ymin=0, ymax=4,
    axis lines=left,
    xtick=\empty,
    ytick=\empty,
    title={\small Bremsen}
]
\addplot[thick, black] coordinates {(0,3.5) (5,0)};
\end{axis}
\end{tikzpicture}
\end{center}

\textbf{Weg aus $v(t)$-Diagramm:} Fläche unter dem Graphen berechnen!
\begin{itemize}[nosep]
\item Rechteck: $s = v \cdot t$
\item Dreieck: $s = \dfrac{1}{2} \cdot t \cdot v$
\item Trapez: $s = \dfrac{1}{2} \cdot (v_1 + v_2) \cdot t$
\end{itemize}

\abschnitt{Luftwiderstand und Grenzgeschwindigkeit (qualitativ)}

\textit{Keine Berechnungen -- nur Verständnis der Zusammenhänge!}

\formelbox{Luftwiderstand}{
\begin{center}
$F_W \sim v^2 \cdot A$
\end{center}
\begin{itemize}[nosep]
\item Je schneller ($v$ größer) $\rightarrow$ Luftwiderstand \textbf{viel} größer (quadratisch!)
\item Je größere Fläche ($A$ größer) $\rightarrow$ Luftwiderstand größer
\end{itemize}
}

\formelbox{Grenzgeschwindigkeit}{
Wenn Luftwiderstand = Gewichtskraft ($F_W = F_G$), fällt der Körper mit \textbf{konstanter Grenzgeschwindigkeit} weiter.

\vspace{0.2cm}
\begin{tabular}{@{}ll@{}}
Große Fläche $A$ & $\rightarrow$ kleine Grenzgeschwindigkeit (Fallschirm!) \\
Große Masse $m$ & $\rightarrow$ große Grenzgeschwindigkeit \\
Dünne Luft (Höhe) & $\rightarrow$ große Grenzgeschwindigkeit \\
\end{tabular}
}

\textbf{Wichtige Fakten für die Klausur:}
\begin{itemize}[nosep]
\item \textbf{Im Vakuum:} Kein Luftwiderstand $\rightarrow$ alle Körper fallen gleich schnell (Feder = Münze)
\item \textbf{In Luft:} Leichte Objekte mit großer Fläche (Feder) erreichen schnell ihre niedrige Grenzgeschwindigkeit
\item \textbf{Fallschirm:} Große Fläche $\rightarrow$ großer Luftwiderstand $\rightarrow$ kleine Grenzgeschwindigkeit
\item \textbf{Regentropfen:} Ohne Luftwiderstand würden sie mit $\sim$200 m/s aufschlagen -- tatsächlich viel langsamer wegen Grenzgeschwindigkeit
\item \textbf{Formel $s = \frac{1}{2}gt^2$:} Enthält keine Masse! Deshalb fallen im Vakuum alle Körper gleich.
\end{itemize}

\newpage

% === TEIL 2: MUSTERAUFGABEN ===
\section{Musteraufgaben mit Lösungen}

\abschnitt{Einheitenumrechnung}

\begin{beispiel}{Geschwindigkeit umrechnen}
Ein Auto fährt mit 72 km/h. Berechne die Geschwindigkeit in m/s.

\textbf{Lösung:}
\begin{align*}
v &= \frac{72\,\mathrm{km/h}}{3{,}6} = 20\,\mathrm{m/s}
\end{align*}
\end{beispiel}

\abschnitt{Freier Fall}

\begin{beispiel}{Fallzeit berechnen}
Ein Stein fällt von einer Brücke (Höhe 20 m). Berechne die Fallzeit.

\textbf{Lösung:}

Formel: $s = \dfrac{1}{2} g \cdot t^2$

Umstellen nach $t$: $t = \sqrt{\dfrac{2s}{g}}$

Einsetzen: $t = \sqrt{\dfrac{2 \cdot 20\,\mathrm{m}}{10\,\mathrm{m/s^2}}} = \sqrt{4\,\mathrm{s^2}} = 2\,\mathrm{s}$
\end{beispiel}

\begin{beispiel}{Aufprallgeschwindigkeit berechnen}
Mit welcher Geschwindigkeit trifft der Stein aus Beispiel 2 auf?

\textbf{Lösung:}

Formel: $v = \sqrt{2 \cdot g \cdot s}$

Einsetzen: $v = \sqrt{2 \cdot 10\,\mathrm{m/s^2} \cdot 20\,\mathrm{m}} = \sqrt{400\,\mathrm{m^2/s^2}} = 20\,\mathrm{m/s}$

\textit{Alternativ:} $v = g \cdot t = 10\,\mathrm{m/s^2} \cdot 2\,\mathrm{s} = 20\,\mathrm{m/s}$
\end{beispiel}

\abschnitt{Bremsvorgang}

\begin{beispiel}{Bremszeit und Bremsweg}
Ein LKW fährt mit 90 km/h und bremst mit $a = -5\,\mathrm{m/s^2}$.

a) Berechne die Bremszeit.\\
b) Berechne den Bremsweg.

\textbf{Lösung:}

Zuerst: $v_0 = 90\,\mathrm{km/h} = 25\,\mathrm{m/s}$

a) Bremszeit:
\begin{align*}
t &= \frac{v_0}{|a|} = \frac{25\,\mathrm{m/s}}{5\,\mathrm{m/s^2}} = 5\,\mathrm{s}
\end{align*}

b) Bremsweg:
\begin{align*}
s &= \frac{v_0^2}{2 \cdot |a|} = \frac{(25\,\mathrm{m/s})^2}{2 \cdot 5\,\mathrm{m/s^2}} = \frac{625\,\mathrm{m^2/s^2}}{10\,\mathrm{m/s^2}} = 62{,}5\,\mathrm{m}
\end{align*}
\end{beispiel}

\newpage

\abschnitt{Senkrechter Wurf}

\begin{beispiel}{Steigzeit und maximale Höhe}
Ein Ball wird mit $v_0 = 20\,\mathrm{m/s}$ senkrecht nach oben geworfen.

a) Berechne die Steigzeit.\\
b) Berechne die maximale Höhe.\\
c) Mit welcher Geschwindigkeit kommt er zurück?

\textbf{Lösung:}

a) Steigzeit (bis $v = 0$):
\begin{align*}
t_{\text{steig}} &= \frac{v_0}{g} = \frac{20\,\mathrm{m/s}}{10\,\mathrm{m/s^2}} = 2\,\mathrm{s}
\end{align*}

b) Maximale Höhe:
\begin{align*}
h_{\text{max}} &= \frac{v_0^2}{2g} = \frac{(20\,\mathrm{m/s})^2}{2 \cdot 10\,\mathrm{m/s^2}} = \frac{400}{20}\,\mathrm{m} = 20\,\mathrm{m}
\end{align*}

c) Rückkehrgeschwindigkeit: $v = 20\,\mathrm{m/s}$ (nach unten)

Begründung: Energieerhaltung -- die potentielle Energie am höchsten Punkt wird vollständig in kinetische Energie umgewandelt.
\end{beispiel}

\abschnitt{Diagrammauswertung}

\begin{beispiel}{Weg aus $v(t)$-Diagramm}
Das Diagramm zeigt die Fahrt eines Autos:

\begin{center}
\begin{tikzpicture}
\begin{axis}[
    width=10cm, height=4cm,
    xlabel={$t$ in s},
    ylabel={$v$ in m/s},
    xmin=0, xmax=15,
    ymin=0, ymax=12,
    grid=major,
    axis lines=left,
    xtick={0,3,6,9,12,15},
    ytick={0,2,4,6,8,10,12}
]
\addplot[fill=black!10, draw=none] coordinates {(0,0) (6,0) (6,6) (0,0)} \closedcycle;
\addplot[fill=black!10, draw=none] coordinates {(6,0) (12,0) (12,6) (6,6)} \closedcycle;
\addplot[fill=black!10, draw=none] coordinates {(12,0) (15,0) (12,6)} \closedcycle;
\addplot[thick, black] coordinates {(0,0) (6,6) (12,6) (15,0)};
\node at (axis cs:3,2) {\small A};
\node at (axis cs:9,3) {\small B};
\node at (axis cs:13.5,2) {\small C};
\end{axis}
\end{tikzpicture}
\end{center}

Berechne den Gesamtweg.

\textbf{Lösung:}

Phase A (Dreieck): $s_A = \dfrac{1}{2} \cdot 6\,\mathrm{s} \cdot 6\,\mathrm{m/s} = 18\,\mathrm{m}$

Phase B (Rechteck): $s_B = 6\,\mathrm{s} \cdot 6\,\mathrm{m/s} = 36\,\mathrm{m}$

Phase C (Dreieck): $s_C = \dfrac{1}{2} \cdot 3\,\mathrm{s} \cdot 6\,\mathrm{m/s} = 9\,\mathrm{m}$

Gesamtweg: $s = 18 + 36 + 9 = 63\,\mathrm{m}$
\end{beispiel}

\newpage

\abschnitt{Weg-Zeit-Gesetz aufstellen}

\begin{beispiel}{Bewegungsgleichung aufstellen und auswerten}
Ein Zug fährt mit konstant $v = 25\,\mathrm{m/s}$. Bei $t = 0$ ist er an Position $s_0 = 50\,\mathrm{m}$.

a) Stelle das Weg-Zeit-Gesetz auf.\\
b) Wo ist der Zug nach 8 s?\\
c) Wann erreicht er Position 300 m?

\textbf{Lösung:}

a) Weg-Zeit-Gesetz:
\begin{align*}
s(t) &= v \cdot t + s_0 = 25\,\mathrm{m/s} \cdot t + 50\,\mathrm{m}
\end{align*}

b) Position nach 8 s:
\begin{align*}
s(8) &= 25 \cdot 8 + 50 = 200 + 50 = 250\,\mathrm{m}
\end{align*}

c) Zeit für 300 m:
\begin{align*}
300 &= 25t + 50\\
250 &= 25t\\
t &= 10\,\mathrm{s}
\end{align*}
\end{beispiel}

\abschnitt{Überholvorgang / Relativbewegung}

\begin{beispiel}{Überholvorgang berechnen}
Ein PKW (72 km/h) fährt 20 m hinter einem LKW (54 km/h).

a) Berechne die Relativgeschwindigkeit.\\
b) Nach welcher Zeit ist der PKW auf gleicher Höhe?

\textbf{Lösung:}

Umrechnen: $v_{\text{PKW}} = 20\,\mathrm{m/s}$, $v_{\text{LKW}} = 15\,\mathrm{m/s}$

a) Relativgeschwindigkeit:
\begin{align*}
v_{\text{rel}} &= v_{\text{PKW}} - v_{\text{LKW}} = 20 - 15 = 5\,\mathrm{m/s}
\end{align*}

b) Zeit zum Einholen:
\begin{align*}
t &= \frac{s}{v_{\text{rel}}} = \frac{20\,\mathrm{m}}{5\,\mathrm{m/s}} = 4\,\mathrm{s}
\end{align*}
\end{beispiel}

\newpage

\abschnitt{Fall mit Anfangsgeschwindigkeit}

\begin{beispiel}{Senkrechter Wurf nach unten}
Ein Ball wird von einem 80 m hohen Turm mit $v_0 = 10\,\mathrm{m/s}$ nach unten geworfen.

a) Berechne die Fallzeit.\\
b) Berechne die Aufprallgeschwindigkeit.

\textbf{Lösung:}

a) Fallzeit mit Weg-Zeit-Gesetz:
\begin{align*}
s &= v_0 \cdot t + \frac{1}{2} g \cdot t^2\\
80 &= 10t + 5t^2\\
5t^2 + 10t - 80 &= 0\\
t^2 + 2t - 16 &= 0
\end{align*}

Mit Mitternachtsformel: $t = \dfrac{-2 + \sqrt{4 + 64}}{2} = \dfrac{-2 + \sqrt{68}}{2} \approx 3{,}1\,\mathrm{s}$

b) Aufprallgeschwindigkeit:
\begin{align*}
v &= v_0 + g \cdot t = 10 + 10 \cdot 3{,}1 = 41{,}1\,\mathrm{m/s}
\end{align*}

\textit{Alternativ mit zeitloser Formel:}
\begin{align*}
v &= \sqrt{v_0^2 + 2gs} = \sqrt{100 + 1600} = \sqrt{1700} \approx 41{,}2\,\mathrm{m/s}
\end{align*}
\end{beispiel}

\newpage

% === TEIL 3: ÜBUNGSAUFGABEN ===
\section{Übungsaufgaben}

\noindent\fbox{\parbox{\dimexpr\textwidth-2\fboxsep-2\fboxrule}{%
\textbf{Hinweis:} Alle Aufgaben sind ohne Taschenrechner lösbar. Verwende $g = 10\,\mathrm{m/s^2}$.
}}

\vspace{0.5cm}

\uebung{Rechne um: a) 108 km/h = \underline{\hspace{2cm}} m/s \quad b) 15 m/s = \underline{\hspace{2cm}} km/h}

\vspace{0.5cm}

\uebung{Ein Ball fällt vom Balkon (Höhe 45 m). Berechne die Fallzeit.}

\vspace{3cm}

\uebung{Ein Motorrad fährt mit 72 km/h und bremst mit $a = -4\,\mathrm{m/s^2}$.\\
a) Berechne die Bremszeit.\\
b) Berechne den Bremsweg.\\
c) Kann der Fahrer vor einem Hindernis in 40 m Entfernung anhalten?}

\vspace{5cm}

\uebung{Ein Spieler wirft einen Ball mit $v_0 = 10\,\mathrm{m/s}$ senkrecht nach oben.\\
a) Berechne die Steigzeit.\\
b) Berechne die maximale Höhe.}

\vspace{4cm}

\uebung{Das $v(t)$-Diagramm zeigt eine Autofahrt. Berechne:
a) Die Beschleunigung in Phase A.\\
b) Den Gesamtweg.}

\begin{center}
\begin{tikzpicture}
\begin{axis}[
    width=10cm, height=4cm,
    xlabel={$t$ in s},
    ylabel={$v$ in m/s},
    xmin=0, xmax=20,
    ymin=0, ymax=15,
    grid=major,
    axis lines=left,
    xtick={0,4,8,12,16,20},
    ytick={0,3,6,9,12,15}
]
\addplot[fill=black!10, draw=none] coordinates {(0,0) (4,0) (4,12) (0,0)} \closedcycle;
\addplot[fill=black!10, draw=none] coordinates {(4,0) (16,0) (16,12) (4,12)} \closedcycle;
\addplot[fill=black!10, draw=none] coordinates {(16,0) (20,0) (16,12)} \closedcycle;
\addplot[thick, black] coordinates {(0,0) (4,12) (16,12) (20,0)};
\node at (axis cs:2,4) {\small A};
\node at (axis cs:10,6) {\small B};
\node at (axis cs:18,4) {\small C};
\end{axis}
\end{tikzpicture}
\end{center}

\vspace{3cm}

\newpage

\uebung{Ein Zug fährt mit konstant 30 m/s. Bei $t = 0$ ist er an Position $s_0 = 100$ m.\\
a) Stelle das Weg-Zeit-Gesetz auf.\\
b) Wo ist der Zug nach 5 s?\\
c) Wann erreicht er die Position 400 m?}

\vspace{5cm}

\uebung{Zwei Straßenbahnen fahren auf parallelen Gleisen. Bahn A fährt mit konstant 15 m/s. Bahn B startet 100 m hinter Bahn A aus dem Stand und beschleunigt mit $a = 2\,\mathrm{m/s^2}$.\\
Nach welcher Zeit hat Bahn B die Bahn A eingeholt?\\
\textit{Tipp: Stelle für beide Bahnen das Weg-Zeit-Gesetz auf und setze gleich.}}

\vspace{5cm}

\uebung{Ein Regentropfen fällt aus einer Wolke in 2000 m Höhe.\\
a) Berechne die theoretische Aufprallgeschwindigkeit (ohne Luftwiderstand).\\
b) Erkläre, warum die tatsächliche Geschwindigkeit viel geringer ist.}

\vspace{4cm}

\newpage

% === TEIL 4: LÖSUNGEN ===
\section{Lösungen zu den Übungsaufgaben}

\noindent\textbf{Übung 1:} a) 30 m/s \quad b) 54 km/h

\vspace{0.3cm}

\noindent\textbf{Übung 2:}\\
$t = \sqrt{\dfrac{2s}{g}} = \sqrt{\dfrac{2 \cdot 45}{10}} = \sqrt{9} = 3\,\mathrm{s}$

\vspace{0.3cm}

\noindent\textbf{Übung 3:}\\
$v_0 = 72 : 3{,}6 = 20\,\mathrm{m/s}$\\
a) $t = \dfrac{v_0}{|a|} = \dfrac{20}{4} = 5\,\mathrm{s}$\\
b) $s = \dfrac{v_0^2}{2|a|} = \dfrac{400}{8} = 50\,\mathrm{m}$\\
c) Nein, Bremsweg (50 m) > Abstand (40 m). Es fehlen 10 m.

\vspace{0.3cm}

\noindent\textbf{Übung 4:}\\
a) $t_{\text{steig}} = \dfrac{v_0}{g} = \dfrac{10}{10} = 1\,\mathrm{s}$\\
b) $h_{\text{max}} = \dfrac{v_0^2}{2g} = \dfrac{100}{20} = 5\,\mathrm{m}$

\vspace{0.3cm}

\noindent\textbf{Übung 5:}\\
a) $a = \dfrac{\Delta v}{\Delta t} = \dfrac{12}{4} = 3\,\mathrm{m/s^2}$\\
b) Phase A: $s_A = \dfrac{1}{2} \cdot 4 \cdot 12 = 24\,\mathrm{m}$\\
\phantom{b)} Phase B: $s_B = 12 \cdot 12 = 144\,\mathrm{m}$\\
\phantom{b)} Phase C: $s_C = \dfrac{1}{2} \cdot 4 \cdot 12 = 24\,\mathrm{m}$\\
\phantom{b)} Gesamt: $s = 24 + 144 + 24 = 192\,\mathrm{m}$

\vspace{0.3cm}

\noindent\textbf{Übung 6:}\\
a) $s(t) = 30t + 100$\\
b) $s(5) = 150 + 100 = 250\,\mathrm{m}$\\
c) $400 = 30t + 100 \Rightarrow t = 10\,\mathrm{s}$

\vspace{0.3cm}

\noindent\textbf{Übung 7:}\\
Bahn A: $s_A(t) = 15t$\\
Bahn B: $s_B(t) = \dfrac{1}{2} \cdot 2 \cdot t^2 - 100 = t^2 - 100$\\
Gleichsetzen: $t^2 - 100 = 15t \Rightarrow t^2 - 15t - 100 = 0$\\
$(t-20)(t+5) = 0 \Rightarrow t = 20\,\mathrm{s}$

\vspace{0.3cm}

\noindent\textbf{Übung 8:}\\
a) $v = \sqrt{2gs} = \sqrt{2 \cdot 10 \cdot 2000} = \sqrt{40000} = 200\,\mathrm{m/s}$\\
b) Der Regentropfen erreicht seine Grenzgeschwindigkeit, wenn Luftwiderstand = Gewichtskraft. Dann fällt er mit konstanter (viel niedrigerer) Geschwindigkeit.

\newpage

% === CHECKLISTE ===
\section{Checkliste zur Klausurvorbereitung}

\noindent\fbox{\parbox{\dimexpr\textwidth-2\fboxsep-2\fboxrule}{%
\textbf{Vor der Klausur prüfen:}
}}

\vspace{0.3cm}

\begin{tabular}{|p{0.5cm}|p{14cm}|}
\hline
$\square$ & Ich kann km/h in m/s umrechnen (und umgekehrt). \\
\hline
$\square$ & Ich kenne die wichtigen Umrechnungen: 36→10, 72→20, 90→25, 108→30. \\
\hline
$\square$ & Ich kann die Fallzeit aus der Höhe berechnen: $t = \sqrt{\dfrac{2h}{g}}$ \\
\hline
$\square$ & Ich kann die Aufprallgeschwindigkeit berechnen: $v = \sqrt{2gh}$ \\
\hline
$\square$ & Ich kann Bremszeit und Bremsweg berechnen. \\
\hline
$\square$ & Ich kann die Steigzeit und maximale Höhe beim senkrechten Wurf berechnen. \\
\hline
$\square$ & Ich weiß, was Steigung und Fläche in $s(t)$- und $v(t)$-Diagrammen bedeuten. \\
\hline
$\square$ & Ich kann den Weg als Fläche unter dem $v(t)$-Graphen berechnen. \\
\hline
$\square$ & Ich kann ein Weg-Zeit-Gesetz aufstellen und auswerten. \\
\hline
$\square$ & Ich kann zwei Bewegungsgleichungen gleichsetzen (Überholaufgaben). \\
\hline
$\square$ & Ich kann erklären, warum im Vakuum alle Körper gleich schnell fallen. \\
\hline
$\square$ & Ich weiß, wie Grenzgeschwindigkeit entsteht ($F_W = F_G$) und wovon sie abhängt. \\
\hline
$\square$ & Ich verwende immer $g = 10\,\mathrm{m/s^2}$ und rechne im Kopf. \\
\hline
$\square$ & Ich schreibe bei jeder Rechnung die Formel, stelle um, setze mit Einheiten ein. \\
\hline
\end{tabular}

\vspace{1cm}

\noindent\fbox{\parbox{\dimexpr\textwidth-2\fboxsep-2\fboxrule}{%
\textbf{Typische Fehler vermeiden:}
\begin{itemize}[nosep]
\item Einheiten vergessen: $v = 20$ statt $v = 20\,\mathrm{m/s}$
\item km/h nicht umrechnen: 72 km/h direkt in Formeln einsetzen
\item $g = 9{,}81$ statt $g = 10$ verwenden
\item Beim Bremsen: Bremsweg mit $a$ statt $|a|$ berechnen
\item Steigung/Fläche im Diagramm verwechseln
\item Bei Überholaufgaben: Startposition vergessen
\end{itemize}
}}

\vspace{1cm}

\begin{center}
\textit{Viel Erfolg bei der Klausur!}
\end{center}

\end{document}
